\section{Struttura della tesi}

Il resto della tesi è organizzato come segue. Il Capitolo~2 presenta lo stato dell’arte sul platooning, sui controllori longitudinali, sulle tecnologie V2X considerate e sugli approcci multi-canale, fino a discutere il meccanismo \emph{SafeSwitch} che costituisce il riferimento diretto del lavoro sperimentale. Il Capitolo~3 descrive invece l’ambiente di sviluppo e simulazione, chiarendo il ruolo di OMNeT++, Veins, PLEXE, SUMO e Simu5G nella toolchain adottata.

Il Capitolo~4 affronta l’importazione della logica del progetto di riferimento, concentrandosi su scenari, metriche, integrazione del PDR, meccanismo \emph{SafeSwitch} e passaggio della componente cellulare da SimuLTE a Simu5G. Il Capitolo~5 presenta e discute i risultati ottenuti, includendo sia il confronto con un fallback \emph{naive} sia l’analisi dello scenario realistico con 5G. Infine, il Capitolo~6 riassume i principali contributi della tesi, ne discute i limiti e delinea alcuni possibili sviluppi futuri.

Prima dell’appendice è inoltre riportata una tabella degli acronimi impiegati nella trattazione, così da rendere più rapida la consultazione delle principali sigle tecniche.

L’Appendice~\ref{ch:setup} dettaglia la preparazione dell’ambiente sperimentale, includendo organizzazione del workspace, build e validazione operativa della toolchain.
